% 示例：金属电子逸出功的测定（深圳大学《大学物理实验（2）》格式）
% 数据为一次真实实验记录；封面个人信息已替换为占位。
% 编译：python analysis.py && python ../../skills/lab-report-tex/scripts/build.py report.tex --thumbs preview
\documentclass{labreport}
\def\ExtKa{0.008075}%  fit_plot.py lgia_long.csv lgIa~sqrtUa
\def\ExtBa{-1.446}%  fit_plot.py lgia_long.csv lgIa~sqrtUa
\def\ExtRa{0.5774}%  fit_plot.py lgia_long.csv lgIa~sqrtUa
\def\ExtUKa{0.0047}%  fit_plot.py lgia_long.csv lgIa~sqrtUa
\def\ExtUBa{0.037}%  fit_plot.py lgia_long.csv lgIa~sqrtUa
\def\ExtKb{0.009286}%  fit_plot.py lgia_long.csv lgIa~sqrtUa
\def\ExtBb{-0.8905}%  fit_plot.py lgia_long.csv lgIa~sqrtUa
\def\ExtRb{0.9421}%  fit_plot.py lgia_long.csv lgIa~sqrtUa
\def\ExtUKb{0.0013}%  fit_plot.py lgia_long.csv lgIa~sqrtUa
\def\ExtUBb{0.011}%  fit_plot.py lgia_long.csv lgIa~sqrtUa
\def\ExtKc{0.008267}%  fit_plot.py lgia_long.csv lgIa~sqrtUa
\def\ExtBc{-0.4068}%  fit_plot.py lgia_long.csv lgIa~sqrtUa
\def\ExtRc{0.9879}%  fit_plot.py lgia_long.csv lgIa~sqrtUa
\def\ExtUKc{5.3e-04}%  fit_plot.py lgia_long.csv lgIa~sqrtUa
\def\ExtUBc{0.0042}%  fit_plot.py lgia_long.csv lgIa~sqrtUa
\def\ExtKd{0.007857}%  fit_plot.py lgia_long.csv lgIa~sqrtUa
\def\ExtBd{0.02215}%  fit_plot.py lgia_long.csv lgIa~sqrtUa
\def\ExtRd{0.9980}%  fit_plot.py lgia_long.csv lgIa~sqrtUa
\def\ExtUKd{2.0e-04}%  fit_plot.py lgia_long.csv lgIa~sqrtUa
\def\ExtUBd{0.0016}%  fit_plot.py lgia_long.csv lgIa~sqrtUa
\def\ExtKe{0.007288}%  fit_plot.py lgia_long.csv lgIa~sqrtUa
\def\ExtBe{0.4115}%  fit_plot.py lgia_long.csv lgIa~sqrtUa
\def\ExtRe{0.9991}%  fit_plot.py lgia_long.csv lgIa~sqrtUa
\def\ExtUKe{1.2e-04}%  fit_plot.py lgia_long.csv lgIa~sqrtUa
\def\ExtUBe{9.7e-04}%  fit_plot.py lgia_long.csv lgIa~sqrtUa
\def\RichK{-2.647}%  fit_plot.py rich_fit.csv y~x
\def\RichB{5.722}%  fit_plot.py rich_fit.csv y~x
\def\RichR{-0.9994}%  fit_plot.py rich_fit.csv y~x
\def\RichUK{0.053}%  fit_plot.py rich_fit.csv y~x
\def\RichUB{0.26}%  fit_plot.py rich_fit.csv y~x
\def\RichKK{-2.647e+04}%  analysis.py
\def\RichUKK{5.3e+02}%  analysis.py
\def\KbLnE{1.984e-04}%  analysis.py
\def\Vraw{5.252}%  analysis.py
\def\uVraw{0.105}%  analysis.py
\def\WVal{5.25}%  analysis.py
\def\WU{0.11}%  analysis.py
\def\WRel{2.0}%  analysis.py
\def\WZero{4.54}%  analysis.py
\def\WErr{15.6}%  analysis.py
\def\IzeroMin{0.036}%  analysis.py
\def\IzeroMax{2.579}%  analysis.py
\def\ExtRmin{0.5774}%  analysis.py
\def\ExtRmax{0.9991}%  analysis.py


\labset{
  school     = {深圳大学},
  cover      = {szu},
  coursename = {大学物理实验（2）},
  coursecode = {1000001-00},
  expname    = {金属电子逸出功的测定},
  college    = {物理与电子科学学院},
  teacher    = {李老师},
  author     = {张三},
  studentid  = {2025000000},
  group      = {1},
  location   = {实验楼A301},
  date       = {2026年10月10日},
  submitdate = {2026年10月12日},
}

\begin{document}
\makecover

\begin{labsection}{实验目的}
1．了解金属热电子发射的基本规律；\par
2．学习用外延法测定不同灯丝温度下的零场热电子发射电流；\par
3．学习用理查逊直线法测定钨的逸出电势和逸出功；\par
4．掌握作图法和最小二乘拟合处理实验数据的方法。
\end{labsection}

\begin{labsection}{实验原理}
\labsub{热电子发射与理查逊-杜西曼公式}
真空二极管的阴极用被测金属钨丝制成，通电加热后阴极温度升高，金属内电子的能量分布改变，动能大于逸出功的电子数目增多，这些电子能够克服表面势垒从金属表面逸出，形成热电子发射。在阳极未加电压时，外电路中也存在热发射电流，称为零场电流。零场电流由理查逊-杜西曼公式确定：
\begin{equation}
  I = AST^2 \exp\left(-\frac{eV}{k_\mathrm{B}T}\right)
  \label{eq:rd}
\end{equation}
式中 $A$ 为与阴极表面化学纯度有关的系数，单位为 \si{A/(m^2.K^2)}；$S$ 为阴极的有效发射面积；$T$ 为阴极绝对温度；$k_\mathrm{B}$ 为玻尔兹曼常数；$e$ 为电子电荷；$V$ 为电子的逸出电势，逸出功 $W=eV$。

将式 \eqref{eq:rd} 两边除以 $T^2$ 再取常用对数，得
\begin{equation}
  \lg\frac{I}{T^2} = \lg(AS) - \frac{eV\lg e}{k_\mathrm{B}}\cdot\frac{1}{T}
  \label{eq:line}
\end{equation}
即 $\lg(I/T^2)$ 与 $1/T$ 成线性关系。由直线斜率 $k=-eV\lg e/k_\mathrm{B}$ 即可求出逸出电势 $V$，这种方法称为理查逊直线法。它避开了难以准确测定的 $A$ 和 $S$。

\labsub{零场电流的测量——外延法}
热电子飞向阳极时形成空间电荷，阻碍后续电子，因此需在阳极加加速电压 $U_a$；但 $U_a$ 又会产生肖脱基效应，使测得的阳极电流 $I_a$ 大于零场电流 $I_0$。可以证明
\begin{equation}
  \lg I_a = \lg I_0 + b\sqrt{U_a}
  \label{eq:schottky}
\end{equation}
以 $\lg I_a$ 为纵坐标、$\sqrt{U_a}$ 为横坐标作图，将直线外延到 $\sqrt{U_a}=0$，截距即为 $\lg I_0$。由式 \eqref{eq:line}，
\begin{equation}
  V = -k\cdot\frac{k_\mathrm{B}\ln 10}{e},\qquad \frac{k_\mathrm{B}\ln 10}{e}=\num{\KbLnE}\,\si{V/K}
  \label{eq:v}
\end{equation}
\end{labsection}

\begin{labsection}{实验仪器}
THQYC-1 型金属电子逸出功实验仪、THQYC-1 型金属电子逸出功测试台（内装真空二极管）。主要仪器及参数如下：
\labtable{主要仪器及参数}[tab:instr]{
\begin{tabularx}{\linewidth}{|>{\centering\arraybackslash}p{4.2cm}|>{\centering\arraybackslash}X|>{\centering\arraybackslash}X|}
\hline 仪器名称 & 主要用途 & 参数 \\ \hline
金属电子逸出功实验仪 & 提供灯丝电流 $I_\mathrm{f}$ 与阳极电压 $U_a$ & $I_\mathrm{f}$ 分度 0.01\,A；$U_a$ 16～121\,V \\ \hline
逸出功测试台 & 安装真空二极管及微安表 & 微安表分度值 0.01\,μA \\ \hline
真空二极管 & 热电子发射源（钨丝阴极） & 阳极、阴极共轴圆柱形 \\ \hline
\end{tabularx}}
\end{labsection}

\begin{labsection}{实验内容与步骤}
1．按要求接好电路，接通电源，仪器预热约 10 分钟。\par
2．灯丝电流 $I_\mathrm{f}$ 在 0.52～0.68\,A 范围内每隔 0.04\,A 取一个值，共 5 个。\par
3．对应每一灯丝电流，阳极电压 $U_a$ 依次取 16、25、36、49、64、81、100、121\,V，记录阳极电流 $I_a$。\par
4．对每组数据作 $\lg I_a$-$\sqrt{U_a}$ 图，用外延法求出各温度下的零场电流 $I_0$。\par
5．作 $\lg(I_0/T^2)$-$1/T$ 图，由斜率求出逸出电势 $V$ 和逸出功 $W$，与公认值比较。\par
\textbf{注意事项：}灯丝电流不可超过额定值；每次改变 $I_\mathrm{f}$ 后须预热几分钟，待灯丝温度稳定后再读数。
\end{labsection}

\begin{labsection}{数据记录和处理}
\textbf{1．数据记录}\par
灯丝电流与温度的对应关系见表 \ref{tab:ift}，各灯丝电流下的阳极电流见表 \ref{tab:ia}（原始记录见最后的原始数据页）。
% 由 csv2tex.py 从 t_ift.csv 生成，改数据请改 CSV 后重新生成
\labtable{灯丝电流与灯丝温度的对应关系}[tab:ift]{
\begin{tabular}{cc}
\toprule
$I_\mathrm{f}$/A & $T$/$10^3$K \\
\midrule
0.52 & 1.93 \\
0.56 & 2.00 \\
0.60 & 2.07 \\
0.64 & 2.14 \\
0.68 & 2.21 \\
\bottomrule
\end{tabular}}

% 由 csv2tex.py 从 t_ia.csv 生成，改数据请改 CSV 后重新生成
\labtable{不同灯丝电流和阳极电压$U_a$/V下的阳极电流$I_a$/μA}[tab:ia]{
\begin{tabular}{ccccccccc}
\toprule
$I_\mathrm{f}$/A & 16 & 25 & 36 & 49 & 64 & 81 & 100 & 121 \\
\midrule
0.52 & 0.04 & 0.04 & 0.04 & 0.04 & 0.04 & 0.04 & 0.04 & 0.05 \\
0.56 & 0.14 & 0.14 & 0.15 & 0.15 & 0.15 & 0.16 & 0.16 & 0.16 \\
0.60 & 0.42 & 0.43 & 0.44 & 0.45 & 0.46 & 0.47 & 0.47 & 0.48 \\
0.64 & 1.13 & 1.15 & 1.17 & 1.20 & 1.22 & 1.24 & 1.26 & 1.28 \\
0.68 & 2.75 & 2.81 & 2.86 & 2.90 & 2.95 & 3.00 & 3.05 & 3.10 \\
\bottomrule
\end{tabular}}


\textbf{2．外延法求零场电流}\par
对表 \ref{tab:ia} 取常用对数得表 \ref{tab:lgia}，对每个灯丝电流作 $\lg I_a$-$\sqrt{U_a}$ 最小二乘拟合（图 \ref{fig:lgia}），结果见表 \ref{tab:ext}。
% 由 csv2tex.py 从 t_lgia.csv 生成，改数据请改 CSV 后重新生成
\labtable{$\lg(I_a/\mu\mathrm{A})$ 与 $\sqrt{U_a/\mathrm{V}}$（首行）}[tab:lgia]{
\begin{tabular}{ccccccccc}
\toprule
$I_\mathrm{f}$/A & 4 & 5 & 6 & 7 & 8 & 9 & 10 & 11 \\
\midrule
0.52 & $-1.3979$ & $-1.3979$ & $-1.3979$ & $-1.3979$ & $-1.3979$ & $-1.3979$ & $-1.3979$ & $-1.3010$ \\
0.56 & $-0.8539$ & $-0.8539$ & $-0.8239$ & $-0.8239$ & $-0.8239$ & $-0.7959$ & $-0.7959$ & $-0.7959$ \\
0.60 & $-0.3768$ & $-0.3665$ & $-0.3565$ & $-0.3468$ & $-0.3372$ & $-0.3279$ & $-0.3279$ & $-0.3188$ \\
0.64 & 0.0531 & 0.0607 & 0.0682 & 0.0792 & 0.0864 & 0.0934 & 0.1004 & 0.1072 \\
0.68 & 0.4393 & 0.4487 & 0.4564 & 0.4624 & 0.4698 & 0.4771 & 0.4843 & 0.4914 \\
\bottomrule
\end{tabular}}

\labfig[0.78\linewidth]{fig/lgia.pdf}{$\lg I_a$-$\sqrt{U_a}$ 关系及外延直线}[fig:lgia]
% 由 csv2tex.py 从 t_ext.csv 生成，改数据请改 CSV 后重新生成
\labtable{外延法拟合结果（$\lg I_a=k\sqrt{U_a}+\lg I_0$）}[tab:ext]{
\begin{tabular}{ccccc}
\toprule
$I_\mathrm{f}$/A & 斜率 & $\lg(I_0/\mu\mathrm{A})$ & $r$ & $I_0$/μA \\
\midrule
0.52 & 0.0081 & $-1.4464$ & 0.5774 & 0.036 \\
0.56 & 0.0093 & $-0.8905$ & 0.9421 & 0.129 \\
0.60 & 0.0083 & $-0.4068$ & 0.9879 & 0.392 \\
0.64 & 0.0079 & 0.0221 & 0.9980 & 1.052 \\
0.68 & 0.0073 & 0.4115 & 0.9991 & 2.579 \\
\bottomrule
\end{tabular}}

示例：$I_\mathrm{f}=0.60$\,A 时，拟合得斜率 $k=\num{\ExtKc}$、截距 $\lg I_0=\num{\ExtBc}$，故零场电流 $I_0=10^{\ExtBc}\,\text{μA}=0.392$\,μA。

\textbf{3．理查逊直线法求逸出功}\par
由表 \ref{tab:ift} 的温度和表 \ref{tab:ext} 的零场电流计算 $\lg(I_0/T^2)$ 与 $1/T$（表 \ref{tab:rich}），作图如图 \ref{fig:rich}。
% 由 csv2tex.py 从 rich.csv 生成，改数据请改 CSV 后重新生成
\labtable{不同灯丝温度下的 $\lg(I_0/T^2)$ 与 $1/T$}[tab:rich]{
\begin{tabular}{ccccc}
\toprule
$I_\mathrm{f}$/A & $T$/$10^3$K & $\lg(I_0/\mu\mathrm{A})$ & $\lg(I_0/T^2)$ & $(1/T)$/$10^{-4}$K$^{-1}$ \\
\midrule
0.52 & 1.93 & $-1.4464$ & $-8.0175$ & 5.1813 \\
0.56 & 2.00 & $-0.8905$ & $-7.4926$ & 5.0000 \\
0.60 & 2.07 & $-0.4068$ & $-7.0387$ & 4.8309 \\
0.64 & 2.14 & 0.0221 & $-6.6387$ & 4.6729 \\
0.68 & 2.21 & 0.4115 & $-6.2773$ & 4.5249 \\
\bottomrule
\end{tabular}}

\labfig[0.7\linewidth]{fig/rich.pdf}{$\lg(I_0/T^2)$-$1/T$ 关系曲线}[fig:rich]
最小二乘拟合得斜率 $k=\num{\RichKK}$\,K，相关系数 $r=\num{\RichR}$，斜率的标准不确定度 $u(k)=\num{\RichUKK}$\,K。

\textbf{4．逸出功及其不确定度}\par
由式 \eqref{eq:v}，
\[
  V = -k\cdot\frac{k_\mathrm{B}\ln 10}{e} = -(\num{\RichKK})\times\num{\KbLnE} = \Vraw\ \mathrm{V}
\]
只计拟合斜率的统计不确定度：$u(V)=u(k)\cdot k_\mathrm{B}\ln10/e=\uVraw$\,V。按有效数字规则修约，
\[
  W = (\WVal \pm \WU)\ \mathrm{eV}\quad(P=0.683),\qquad u(W)/W=\WRel\%
\]
与钨的逸出功公认值 $W_0=\WZero$\,eV 比较，相对误差
$E=|W-W_0|/W_0=|\WVal-\WZero|/\WZero=\WErr\%$。
\end{labsection}

\begin{labsection}{结果陈述及总结}
1．由表 \ref{tab:ext} 及图 \ref{fig:lgia}，各灯丝电流下 $\lg I_a$ 随 $\sqrt{U_a}$ 线性增大（$I_\mathrm{f}=0.68$\,A 时 $r=\ExtRmax$），符合肖脱基效应；外延得到的零场电流由 \IzeroMin\,μA 增大到 \IzeroMax\,μA，随灯丝温度升高而增大。\par
2．由图 \ref{fig:rich}，$\lg(I_0/T^2)$ 与 $1/T$ 线性良好（$r=\num{\RichR}$），符合理查逊-杜西曼公式。\par
3．测得钨的逸出功 $W=(\WVal\pm\WU)$\,eV，与公认值 \WZero\,eV 相比相对误差 \WErr\%。\par
4．结果偏大的可能原因：① $I_\mathrm{f}=0.52$\,A 时 $I_a$ 仅 0.04\,μA 量级，接近微安表分辨率 0.01\,μA，该组线性较差（$r=\ExtRmin$），外延所得 $I_0$ 不确定度大；② 灯丝温度按仪器标定的 $I_\mathrm{f}$-$T$ 关系读取，标定本身有系统误差；③ 外延公式为近似关系，未计接触电位差。
\end{labsection}

\begin{labsection}{思考题}
\question{理查逊直线法有何优点？}
把逸出功归结为 $\lg(I/T^2)$-$1/T$ 直线的斜率，只需测量电流和温度两个量，避免了直接测定与阴极表面状况有关的系数 $A$ 和有效发射面积 $S$——这两个量在实验中很难准确测定，因此大大简化了测量并提高了精度。

\question{灯丝电流为何要保持稳定？每次改变 $I_\mathrm{f}$ 后为何要预热几分钟再测量？}
灯丝温度由灯丝电流决定，而由式 \eqref{eq:rd}，发射电流对温度极其敏感，温度稍有变化 $I$ 就显著改变。灯丝电流不稳，数据点就会偏离直线。改变 $I_\mathrm{f}$ 后灯丝需经热传导才能达到新的平衡温度，未稳定就读数，实际温度与标定的 $I_\mathrm{f}$-$T$ 关系不符，测得的数据不可靠。
\end{labsection}

\labcomment
\labgrading{预习/20,操作及记录/40,数据处理与结果陈述/30,思考题/10}

\labphotos{原始数据}{raw/raw1.jpg}
\end{document}
